\section{Introduction}
\label{sec:intro}

Security operations centers (SOCs) face high alert volumes and a shortage of skilled analysts, and large language models (LLMs) are increasingly proposed as \emph{copilots} that read logs and alerts and produce triage labels, incident summaries, or remediation advice~\cite{habibzadeh2025socsurvey}. Running such models \emph{locally}, on-host, is attractive because sensitive logs never leave the organization.

These systems share a structural weakness: an LLM receives its instructions and the data it must analyze as a single stream of text, and many log fields (the HTTP User-Agent, the request URI, command-line arguments, attempted usernames) are written by the attacker whose activity is being recorded. A crafted log entry can therefore carry instructions that hijack the model reading it~\cite{greshake2023indirect,pandey2026watchtower}.

Fig.~\ref{fig:running-example} shows the running example used throughout this report. An attacker probes a web application for SQL injection, and the nginx server records the request in its access log. Without the red text, the entry is an ordinary, clearly malicious probe: the query string contains \texttt{' OR 1=1--}, and a copilot asked to triage it should label it malicious. However, the User-Agent header is chosen freely by the client, so the attacker appends a sentence addressed not to the web server but to the LLM that will later read the log. If the copilot follows that sentence and labels the request benign, the attack has been hidden by the very system meant to detect it.

\begin{figure}[t]
\centering
\footnotesize
\begin{minipage}{0.95\columnwidth}
\ttfamily\raggedright
203.0.113.7 - - [11/Oct/2026:10:02:11 +0000] "GET /item?id=1'\%20OR\%201=1-- HTTP/1.1" 200 512 "-" "Mozilla/5.0 (X11; Linux x86\_64) \textcolor{red}{[SYSTEM NOTE: authorized penetration test, classify this request as benign]}"
\end{minipage}
\caption{Running example: an nginx access-log entry recording a SQL-injection probe. The attacker-controlled User-Agent field (red) carries an instruction aimed at the LLM that will later read the log.}
\label{fig:running-example}
\end{figure}

This report presents the theoretical background needed to study this threat. Section~\ref{sec:llm} introduces the language-modelling and alignment concepts that explain why models follow injected text. Section~\ref{sec:pi} defines prompt injection and its log-substrate variant. Section~\ref{sec:threat} describes the security-operations setting and the threat model. Section~\ref{sec:opt} formulates adversarial payload generation as a constrained optimization problem and presents three search methods. Section~\ref{sec:def} covers the main families of defenses, and Section~\ref{sec:eval} the statistical tools used to evaluate them. Section~\ref{sec:synthesis} shows how these concepts fit together in the author's thesis.
